% Môn: Vật lý
% Lớp: 11
% Chương: Chương III. Điện trường
% Bài: L11C3  Bài 21. Tụ điện
% Số câu: 10
% App đọc file này trực tiếp — sửa rồi Commit trên GitHub, bấm Đồng bộ GitHub .tex

% L11C3  Bài 21. Tụ điện

% ===== Câu 1 =====
% ID: L11C3B21-01-TL
% Mức: TH
\dangbt{Chưa có dạng}
\begin{bt}
% ID: L11C3B21-01-TL
% Mức: TH
Tính năng lượng được giải phóng (hay công phóng điện) khi ta ghép nối hai tụ điện trong Câu 17 theo cách nối dây dương của tụ điện này với dây âm của tụ điện kia. $\nguon{ly11 kntt.pdf}$
\loigiai{
Khi ta ghép nối hai tụ điện trong Bài 21.17 theo cách nối dây dương của tụ điện này với dây âm của tụ điện kia thì điện tích của bộ tụ điện sau khi ghép nối bằng:

Q_(sau) = |Q_(a)-Q_(b)| = 20 $\cdot10^{-6}C$

Năng lượng trên bộ tụ điện sau khi ghép bằng:

$\text{U}_{\text{sau}} = \dfrac{\text{Q}_{\text{sau}}}{\text{C}_{\text{sau}}} = \dfrac{\text{C}_{\text{a}}\text{U}_{\text{a}} - \text{C}_{\text{b}}\text{U}_{\text{b}}}{\text{C}_{\text{a}} + \text{C}_{\text{b}}} = \dfrac{40}{7}\text{V}$

Năng lượng trên bộ tụ điện sau khi ghép nối bằng: $W_{\text{sau}} = \dfrac{C_{\text{sau}}U_{\text{sau}}^{2}}{2} = 0,000057\text{J}$.

Năng lượng đã giải phóng trong quá trình ghép nối là:

$A$ = $(W_(a/tr)+W_(b/tr))$ - W_(sau) = 0, $020743\,\mathrm{J}$
}
\end{bt}

% ===== Câu 2 =====
% ID: L11C3B21-01-TN
% Mức: NB
\begin{ex}
% ID: L11C3B21-01-TN
% Mức: NB
Đại lượng đặc trưng cho khả năng tích điện của tụ điện là: $\nguon{ly11 kntt.pdf}$
\choice
{\True điện dung $C$}
{điện tích $Q$}
{khoảng cách d giữa hai bản tụ.}
{cường độ điện trường.}
\loigiai{
Chọn A.

Đại lượng đặc trưng cho khả năng tích điện của tụ điện là điện dung C.
}
\end{ex}

% ===== Câu 3 =====
% ID: L11C3B21-02-TL
% Mức: TH
\begin{ex}
% ID: L11C3B21-02-TL
% Mức: TH
Một máy hàn bu - lông dùng hiệu điện thế $220\,\mathrm{V}$ không đổi có bộ tụ điện với điện dung $C = 0$, $09\,\mathrm{F}$. a) Tính năng lượng mà bộ tụ điện của máy hàn trên có thể tích được. b) Máy hàn trên có thể phóng điện giải phóng hoàn toàn năng lượng mà bộ tụ điện đã tích được trong khoảng thời gian từ $0,2\,\mathrm{s}$ đến $1\,\mathrm{s}$. Hãy tính công suất phóng điện tối đa của máy hàn đó. CHƯƠNG 4. DÒNG ĐIỆ $N$ KHÔNG ĐỔI. MẠCH ĐIỆ $N$ BÀ̀ $I$ 22. CƯỜNG ĐỘ DÒNG ĐIỆ $N\nguon{ly11 kntt.pdf}$
\choiceTF

\loigiai{
a) Năng lượng mà bộ tụ điện của máy hàn trên có thể tính được là: $W = \dfrac{CU^{2}}{2} = 2178\text{J}$

b) Công suất phóng điện tối đa của máy hàn đó là: $\text{P} = \dfrac{W}{t_{\min}} =$ 10890 $\mathrm{W}$.
}
\end{ex}

% ===== Câu 4 =====
% ID: L11C3B21-02-TN
% Mức: NB
\begin{ex}
% ID: L11C3B21-02-TN
% Mức: NB
Khi trong phòng thi nghiệm chil có một số tụ điện giống nhau với cùng điện dung $C$, muốn thiết kế một bộ tụ điện có điện dung nhỏ hơn $C$ thi: $\nguon{ly11 kntt.pdf}$
\choice
{chắc chắn phải ghép song song các tự điện.}
{\True chắc chắn phải ghép nối tiếp các tụ điện.}
{chắc chắn phải kết hợp cả ghép song song và nối tiếp.}
{không thể thiết kế được bộ tự điện như vậy.}
\loigiai{
Chọn B.

Để bộ tụ có điện dung nhỏ hơn thì phải ghép nối tiếp các tụ điện với nhau.
}
\end{ex}

% ===== Câu 5 =====
% ID: L11C3B21-03-TN
% Mức: NB
\begin{ex}
% ID: L11C3B21-03-TN
% Mức: NB
Hai tụ điện có điện dung lần lượt $C_1 = 1\muF$, $C_2 = 3\muF ghép nối tiếp$. Mắc bộ tụ điện đó vào hai cực của nguồn điện có hiệu điện thế $U = 40\$, $\mathrm{V}$. Điện tích của các tụ điện là: $\nguon{ly11 kntt.pdf}$
\choice
{$Q_1 = 40 \cdot  10^{-6} C và Q_2 = 120 \cdot  10^{-6} C$.}
{\True $Q_1 = Q_2 = 30\cdot 10^{-6}$}
{C. $Q_1 = 7$,5 $\cdot$  10^{-6} $C$ và $Q_2 = 22$,5 $\cdot$  10^{-6} C.}
{$Q_1 = Q_2 = 160 \cdot  10^{-6} C$.}
\loigiai{
Chọn B.

Điện dung bộ tụ: $C = \dfrac{C_{1}C_{2}}{C_{1} + C_{2}} = \dfrac{1.3}{1 + 3} = 0,75\mu F$ Điện tích các tụ điện: Q₁ = Q₂ = $Q$ = CU = 0, 75.40 = 30μ $C$
}
\end{ex}

% ===== Câu 6 =====
% ID: L11C3B21-04-TN
% Mức: NB
\begin{ex}
% ID: L11C3B21-04-TN
% Mức: NB
Năng lượng của điện trường trong một tụ điện đa̋ tích được điện tích q không phụ thuộc vào $\nguon{ly11 kntt.pdf}$
\choice
{điện tích mà tụ điện tích được.}
{hiệu điện thế giữa hai bản tụ điện.}
{\True thời gian đã thực hiện để tích điện cho tụ điện.}
{điện dung của tụ điện.}
\loigiai{
Chọn C.

$\text{W} = \dfrac{1}{2}CU^{2} = \dfrac{Q^{2}}{$ 2 $\mathrm{C}$ } = \dfrac{1}{2}QU $. Năng lượng của điện trường trong một tụ điện đã tích được điện tích q không phụ thuộc vào thời gian đã thực hiện để tích điện cho tụ điện$
}
\end{ex}

% ===== Câu 7 =====
% ID: L11C3B21-05-TN
% Mức: NB
\begin{ex}
% ID: L11C3B21-05-TN
% Mức: NB
Năng lượng của tụ điện bằng $\nguon{ly11 kntt.pdf}$
\choice
{\True công đễ tích điện cho tụ điện.}
{điện thế của các điện tích trên các bản tụ điện.}
{tổng điện thế của các bản tụ điện.}
{khả năng tích điện của tụ điện.}
\loigiai{
Chọn A.

Năng lượng của tụ điện bằng công để tích điện cho tụ điện.
}
\end{ex}

% ===== Câu 8 =====
% ID: L11C3B21-06-TN
% Mức: NB
\begin{ex}
% ID: L11C3B21-06-TN
% Mức: NB
Một tụ điện có điện tích bằng $Q$ và ngắt khỏi nguồn, nếu tăng khoảng cách giữa hai bản tụ điện thì $\nguon{ly11 kntt.pdf}$
\choice
{năng lượng của tụ điện giảm.}
{năng lượng của tụ điện tăng lên do ta đã cung cấp một công làm tăng thế năng của các điện tích.}
{\True năng lượng của tụ điện không thay đổi.}
{năng lượng của tụ điện tăng lên rồi mới giảm.}
\loigiai{
Chọn C.

Năng lượng của tụ không phụ thuộc vào khoảng cách giữa hai bản tụ điện.
}
\end{ex}

% ===== Câu 9 =====
% ID: L11C3B21-07-TN
% Mức: NB
\begin{ex}
% ID: L11C3B21-07-TN
% Mức: NB
Công dụng nào sau đây của một thiết bị không liên quan tới tụ điện? $\nguon{ly11 kntt.pdf}$
\choice
{Tích trữ năng lượng và cung cấp năng lượng.}
{Lưu trữ điện tích.}
{Lọc dòng điện một chiều.}
{\True Cung cấp nhiệt năng ở bàn là, máy sấy,...}
\loigiai{
Chọn D.

Nhiệt lượng toả ra ở tụ điện rất nhỏ.
}
\end{ex}

% ===== Câu 10 =====
% ID: L11C3B21-08-TN
% Mức: NB
\begin{ex}
% ID: L11C3B21-08-TN
% Mức: NB
Hai tụ điện có điện dung lần lượt $C_1 = 2\muF$, $C_2 = 3\muF ghép song song$. Mắc bộ tụ điện đó vào hai cực của nguồn điện có hiệu điện thế $U = 60\$, $\mathrm{V}$. Điện tích của các tụ điện là: $\nguon{ly11 kntt.pdf}$
\choice
{\True $Q_1 = 120 \cdot  10^{-6} C và Q_2 = 180 \cdot  10^{-6} C$.}
{$Q_1 = Q_2 = 72 \cdot  10^{-6}$}
{C. $Q_1 = 3 \cdot  10^{-6} C và Q_2 = 2 \cdot  10^{-6} C$.}
{$Q_1 = Q_2 = 300 \cdot  10^{-6} C$.}
\loigiai{
Chọn A.

Điện tích của mỗi tụ điện: Q₁ = C₁ $U= 120\cdot10^{-6}C$ và Q₂ = C₂ $U$ = 180 $\cdot$ 10^{-6} $C$
}
\end{ex}
